\scanpage{126}
\begin{center}- \mspage{119}\end{center}\margin{IX}

Comme extension panachée $X_\eta$ de $E_\eta$ par $F_\eta$, on prendra
\[
(9.6.2)\qquad X_\eta=T_\ell(A_K)=X_\eta^0\supset X_\eta^1=T_\ell(A_K)^f\supset X_\eta^2=T_\ell(A_K)^t\supset0.
\]
L'isomorphisme d'extensions de $R_\eta$ par $P_\eta$
\[
 X_\eta^1\simeq F_\eta
\]
est clair. L'isomorphisme d'extensions de $Q$ par $R$ :
\[
 X_\eta/X_\eta^2\simeq F_\eta=D(F'_\eta)
\]
est clair grâce au théorème d'orthogonalité 5.2.

Nous disposons donc des données requises pour définir un accouplement (9.5.6), qui ici (comme $Q^*=\underline M'_\ell$) prend la forme (9.1.2). Le fait que cette définition soit compatible avec celle de 9.2 résulte de 9.5.8. Les considérations de 5.3 montrent d'autre part, sous réserve de la validité du théorème de TATE [33] (donc en tous cas dans le cas où car $K=0$) que l'accouplement en question peut se décrire en termes du seul $\ell$-pro-groupe de Barsotti-Tate $X_\eta=T_\ell(A_K)$ sur $S$ (puisque sa filtration en trois crans, et les extensions $F$ et $E$ qui prolongent $X_\eta^1$ et $X_\eta/X_\eta^2$, peuvent se définir canoniquement en termes du seul pro-$\ell$-groupe $X=T_\ell(A_K)$.
\figure{\begin{tikzpicture}[x=1cm,y=1cm]\draw (0,0)--(.12,.1);\draw (0,.1)--(.12,0);\end{tikzpicture}}

\uline{Remarque} 9.7. Ne supposons plus nécessairement que $A_K$ ait réduction semi-stable sur le trait hensélien $S$, et soit $K'$ une extension galoisienne finie de $K$ telle que $A_{K'}$ ait réduction semi-stable sur le normalisé $S'$ de $S$ dans $K'$ (3.6). Avec les notations de (8.1.8), on définit alors, par descente galoisienne à partir de l'accouplement de monodromie pour $A_{K'}$,

\margin{\unsure{437}}\margin{437}\margin{125}
